\begin{table}[h]
\centering
\caption{Determinants of Highway Placement - Discretionary Sample}
\label{tab:discretionary_sample_results}
\begin{threeparttable}
\begin{tabular*}{1.0\textwidth}{@{\extracolsep{\fill}}l*{2}{r}}
\toprule
 & Discretionary Sample & Discretionary Sample with Interacted CNN Logit \\
\midrule
Residential & \makecell[tr]{-0.013{*} \\ (0.006)} & \makecell[tr]{-0.200{**} \\ (0.066)} \\
Black & \makecell[tr]{0.001 \\ (0.010)} & \makecell[tr]{0.053 \\ (0.119)} \\
Residential x Black & \makecell[tr]{-0.015 \\ (0.011)} & \makecell[tr]{-0.139 \\ (0.124)} \\
CNN Logit & \makecell[tr]{0.033{***} \\ (0.004)} & \makecell[tr]{0.061{***} \\ (0.012)} \\
Residential x CNN Logit &  & \makecell[tr]{-0.038{**} \\ (0.012)} \\
Black x CNN Logit &  & \makecell[tr]{0.010 \\ (0.022)} \\
Residential x Black x CNN Logit &  & \makecell[tr]{-0.025 \\ (0.023)} \\
R-squared & 0.076 & 0.096 \\
Observations & 5301 & 5301 \\
\bottomrule
\end{tabular*}
\begin{tablenotes}[flushleft]
\footnotesize
\item This table contains estimates of the impact of residential zoning and majority-Black status on the likelihood of highway placement using a linear probability model on a subset of grid squares designated as discretionary. The discretionary sample excludes grid squares that are intersected by a highway in 1940 or are directly adjacent to a highway in 1940. Standard errors are reported in parenthesis and estimated using a bootstrap procedure with 1,000 draws. The model includes controls for housing, geographic, and demographic characteristics, as well as city fixed effects and a measure of geographic suitability for highway placement, estimated as a logit by a Convolutional Neural Network (CNN). Column 2 also includes the interaction between the CNN logit and the variables of interest.* p<0.10, ** p<0.05, *** p<0.01
\end{tablenotes}
\end{threeparttable}
\end{table}
